\section{Conclusão}
\label{sec:conclusion}

Este artigo investigou o preenchimento de lacunas em imagens multiespectrais de satélite a partir de uma análise orientada por entropia local e organizada em dois blocos metodológicos complementares.
O primeiro bloco examinou métodos clássicos em cenário multissensor, enfatizando qualidade, robustez e custo computacional.
O segundo bloco examinou modelos de aprendizado profundo no domínio Sentinel-2, com foco na comparação interna entre arquiteturas.
Essa separação permitiu formular conclusões mais precisas e evitar generalizações indevidas entre contextos experimentais distintos.

No bloco clássico, os resultados mostraram que não existe um método universalmente dominante.
Kriging, RBF, spline e exemplar\_based concentram o melhor desempenho em várias métricas, mas exibem compensações relevantes entre fidelidade absoluta, robustez ao ruído e tempo de inferência.
Kriging destacou-se como solução particularmente robusta em cenários mais degradados.
RBF e spline mostraram desempenho muito forte em condições menos adversas.
Exemplar\_based destacou-se como compromisso prático de alto interesse, combinando qualidade elevada com custo menos extremo.

No bloco de aprendizado profundo, a U-Net apresentou a melhor qualidade absoluta em todo o domínio avaliado, sugerindo que a propagação multiescala de contexto com conexões de atalho continua sendo uma estratégia muito eficaz para reconstrução de lacunas.
Entretanto, a análise de robustez mostrou que liderança em qualidade não implica estabilidade relativa máxima.
GAN e ViT perderam menos em vários cenários quando a comparação passou a enfatizar degradação percentual ao ruído.
AE e VAE mantiveram papel importante como referências de eficiência e simplicidade, embora abaixo das arquiteturas mais expressivas em fidelidade final.

Do ponto de vista conceitual, o trabalho mostrou que a entropia local é uma variável útil para estruturar a análise do problema, mas sua interpretação depende do método, da métrica e da escala considerada.
Nos métodos clássicos, a relação entre entropia e qualidade é heterogênea e muitas vezes não monotônica.
Nos modelos de aprendizado profundo, a direção do efeito é mais uniforme, embora a magnitude varie entre arquiteturas.
Assim, a entropia funciona menos como marcador universal de dificuldade e mais como moderador do comportamento relativo dos algoritmos.

Em termos práticos, o estudo sugere que a escolha metodológica em preenchimento de lacunas deve ser guiada pelo cenário de uso.
Quando a prioridade é análise multicritério de métodos interpretáveis em contexto multissensor, a família clássica continua oferecendo alternativas competitivas e não pode ser descartada de forma simplista.
Quando a prioridade é qualidade absoluta no domínio Sentinel-2, a U-Net surge como a opção mais forte entre os modelos avaliados.
Quando robustez relativa ou eficiência marginal passam a ser os critérios centrais, outras arquiteturas e outros métodos ganham relevância.

Como continuidade natural, três direções se destacam.
Primeiro, ampliar a avaliação clássica para conjuntos consolidados ainda mais extensos, preservando equilíbrio entre sensores.
Segundo, levar o bloco de aprendizado profundo para cenários multissensor e testar a estabilidade da hierarquia arquitetural observada aqui.
Terceiro, conectar qualidade de reconstrução a tarefas posteriores, verificando em que medida ganhos em PSNR, SSIM ou coerência espectral se traduzem em benefícios reais para pipelines analíticos de sensoriamento remoto.
